% !TEX root = CoeneDylanBP.tex
%%=============================================================================
%% Inleiding
%%=============================================================================

\chapter{\IfLanguageName{dutch}{Inleiding}{Introduction}}%
\label{ch:inleiding}

Trampolinespringen is een technische en complexe sport, waarbij atleten beoordeeld worden op basis van verschillende componenten: moeilijkheidsgraad (D), uitvoering (E), horizontale verplaatsing (HD) en vluchttijd (ToF). De eindscore van een routine, die in competitieverband uit 10 sprongen bestaat, wordt berekend met de formule: $E + HD + D + TOF - P$. Hierbij staat $P$ voor de eventuele aftrek van strafpunten (\textit{penalties}), die door de jury worden toegekend bij specifieke overtredingen of onderbrekingen \autocite{FTTC2024}.

In offici{\"e}le competities is het HDTS Trampoline Measurement Device de standaard voor het objectief meten van HD en ToF. Dit systeem werkt met vier krachtplaten onder de poten van de trampoline om via krachtverhoudingen de exacte landingspositie en timing te bepalen \autocite{Ferger2017,Ferger2019}.

\begin{figure}[ht]
    \centering
    \includegraphics[width=0.8\textwidth]{hdts-trampolinemeasurementsystem}
    \caption[HDTS Trampoline Measurement Device]{De vier krachtplaten van het HDTS Trampoline Measurement Device, gemonteerd onder de poten van de trampoline. Bron: \autocite{Eurotramp2025}.}
    \label{fig:hdts-system}
\end{figure}

\section{\IfLanguageName{dutch}{Probleemstelling}{Problem Statement}}%
\label{sec:probleemstelling}
Ondanks de effectiviteit van de huidige technologie, vormt de hoge kostprijs van het HDTS-systeem een hoge drempel voor veel turnclubs. De prijs van het systeem ligt tussen \euro{}5.000 en \euro{}10.000. De financi{\"e}le druk op de Vlaamse sportsector beperkt de investeringsruimte voor dergelijke apparatuur aanzienlijk \autocite{Sporza2025}.

Turnclubs met beperkte budgetten kunnen zich hierdoor geen HDTS-systeem veroorloven. Dit cre{\"e}ert een ongelijkheid, aangezien atleten in deze clubs geen toegang hebben tot automatische metingen van HD en ToF tijdens hun trainingen. Zij zijn genoodzaakt om hun prestaties te analyseren via tijdsintensieve, manuele video-analyses, wat ineffici{\"e}nt is.

\section{\IfLanguageName{dutch}{Onderzoeksdoelstelling}{Research objective}}%
\label{sec:onderzoeksdoelstelling}

Het doel van dit onderzoek is om een \textit{proof-of-concept} te ontwikkelen van een kosteneffectief alternatief voor commerci{\"e}le HDTS-systemen op basis van computer vision. Het beoogde resultaat is een werkend prototype met een vaste cameraopstelling die voldoet aan de volgende criteria:
\begin{itemize}
    \item Aantonen dat de pipeline HD-landingszones correct kan onderscheiden op basis van videobeelden alleen, als proof-of-concept voor een camera-gebaseerde aanpak.
    \item Een gemiddelde absolute afwijking (MAE) voor de vluchttijd van minder dan 0,2 seconden, bij gebruik van video met voldoende framerate.
    \item Een systeemkost die significant lager ligt dan die van het commerci{\"e}le HDTS-systeem, als ontwerpdoel voor een praktische opstelling.
\end{itemize}

\section{\IfLanguageName{dutch}{Onderzoeksvraag}{Research question}}%
\label{sec:onderzoeksvraag}

De centrale hoofdvraag van deze bachelorproef luidt: 
\textit{Hoe kan een goedkoop computer vision-systeem worden ontwikkeld dat horizontale verplaatsing en vluchttijd bij trampolinespringen betrouwbaar kan benaderen, als betaalbaar alternatief voor het HDTS-meetsysteem?}

Om deze vraag te beantwoorden, worden onder meer de volgende deelvragen behandeld:
\begin{itemize}
    \item Welke stappen in de pipeline vormen de voornaamste bottleneck voor de nauwkeurigheid, en hoe werken fouten van afzonderlijke componenten door naar het eindresultaat?
    \item Welke computer vision-technieken zijn het meest geschikt voor het detecteren van landingsposities en contactmomenten? % Homography transformation, object detection, pose estimation, ...
    \item Welke pre-processing stappen zijn noodzakelijk om correcte inschattingen van HD en ToF te kunnen maken?
    \item Welke state-of-the-art modellen kunnen bijdragen aan een hogere nauwkeurigheid? % pose estimation, depth estimation, ...
    \item Is het effectiever om contactmoment- en positiedetectie te scheiden, of biedt een end-to-end model betere resultaten?
\end{itemize}

\section{\IfLanguageName{dutch}{Opzet van deze bachelorproef}{Structure of this bachelor thesis}}%
\label{sec:opzet-bachelorproef}

De rest van deze bachelorproef is als volgt opgebouwd:

In Hoofdstuk~\ref{ch:stand-van-zaken} wordt een overzicht gegeven van de stand van zaken binnen het onderzoeksdomein, op basis van een literatuurstudie.

In Hoofdstuk~\ref{ch:methodologie} wordt de methodologie toegelicht: het iteratieve verbeterproces waarbij opeenvolgende experimentele fasen telkens aanleiding gaven tot bijsturing van de pipeline, van de initi{\"e}le verkenning en dataverzameling over hoekpuntdetectie en homografietransformatie tot landingsmomentdetectie, landingspositiebepaling en ToF-schatting.

In Hoofdstuk~\ref{ch:resultaten} worden de resultaten van de pipeline besproken per component en getoetst aan de vooropgestelde criteria.

In Hoofdstuk~\ref{ch:discussie} worden de limitaties van het ontwikkelde systeem besproken en aanbevelingen voor verder onderzoek geformuleerd.

In Hoofdstuk~\ref{ch:conclusie}, tenslotte, wordt de conclusie gegeven en een antwoord geformuleerd op de onderzoeksvragen.